% GENERATED FILE. Edit canonical lessons, then run npm run generate:books.
% canonical-source-hash: fnv1a64:4fe6b7ef56709522
% canonical-lessons: HI-C425-halla, HI-C425-khaderna, HI-C425-gatar, HI-C425-badbu, HI-C425-risna

\chapter{Commotion and Drains}
\label{ch:hi-commotion-and-drains}
\hlchaptermodality{425}

\begin{chapteropening}
\textbf{By the end of this chapter:} \emph{I can describe a commotion in the lane, monkeys chased off, a blocked drain and a leak.}

Complete the last lesson of chapter 425: I can describe a commotion in the lane, monkeys chased off, a blocked drain and a leak.
\end{chapteropening}

\section[\texorpdfstring{\hi{हल्ला} (hallā)}{हल्ला (hallā)}]{\hi{हल्ला} (hallā) — an uproar, a commotion}

\label{lesson:HI-C425-halla}



\begin{quote}
Before the new one: say the Hindi for a hawker, then the Hindi for a betel leaf.
\end{quote}

\subsection*{You'll want to know: \hi{हल्ला}}
\textbf{\hi{हल्ला}} — \emph{hallā} — \textquotedblleft{}an uproar, a commotion\textquotedblright{}.

\begin{quote}
\hi{गली} \hi{में} \hi{अचानक} \hi{हल्ला} \hi{हुआ।}
\end{quote}

\begin{quote}
\emph{galī meṁ acānak hallā huā.}
\end{quote}

\begin{quote}
Suddenly there was an uproar in the lane.
\end{quote}

Masculine.

\begin{cousinweb}
\textbf{\hi{हल्ला}} (\emph{hallā}) goes back to \emph{*halla}, \textquotedblleft{}movement\textquotedblright{}, a form CDIAL reconstructs (CDIAL 14017), which lists Hindi \emph{hallā}, \textquotedblleft{}tumult, noise\textquotedblright{}.
\end{cousinweb}

\begin{grammarlens}[title={What you've built}]
One more word for getting around.
\end{grammarlens}

\subsection*{Your turn}
\begin{itemize}
\raggedright
  \item \emph{Say it:} \emph{hallā}
  \item \emph{Say it:} \emph{hallā}, once more
  \item \emph{Say it:} \emph{galī meṁ acānak hallā huā.} — Suddenly there was an uproar in the lane.
  \item \emph{Recall:} say the Hindi for a hawker, then the Hindi for a betel leaf, then say \emph{hallā} again
\end{itemize}

\begin{tcolorbox}[breakable,colback=teal!4,colframe=teal!35!black,title={Before you move on}]
What does \hi{हल्ला} mean? (\textquotedblleft{}An uproar, a commotion\textquotedblright{}.) Say it once more.
\end{tcolorbox}
\section[\texorpdfstring{\hi{खदेड़ना} (khadeṛnā)}{खदेड़ना (khadeṛnā)}]{\hi{खदेड़ना} (khadeṛnā) — to chase away}

\label{lesson:HI-C425-khaderna}



\begin{quote}
Before the new one: say the Hindi for a betel leaf, then the Hindi for an uproar.
\end{quote}

\subsection*{You'll want to know: \hi{खदेड़ना}}
\textbf{\hi{खदेड़ना}} — \emph{khadeṛnā} — \textquotedblleft{}to chase away\textquotedblright{}.

\begin{quote}
\hi{लोगों} \hi{ने} \hi{बंदरों} \hi{को} \hi{गली} \hi{से} \hi{खदेड़ा।}
\end{quote}

\begin{quote}
\emph{logoṁ ne bandroṁ ko galī se khadeṛā.}
\end{quote}

\begin{quote}
People chased the monkeys out of the lane.
\end{quote}

Past, with \hi{ने} (he or she …): \hi{उसने} \hi{बकरी} \hi{को} \hi{खेत} \hi{से} \hi{खदेड़ा}, \emph{usne bakrī ko khet se khadeṛā}.

\begin{cousinweb}
\textbf{\hi{खदेड़ना}} (\emph{khadeṛnā}) goes back to \emph{*khadd}, \textquotedblleft{}drive away\textquotedblright{}, a form CDIAL reconstructs (CDIAL 3807), which lists Hindi \emph{khadeṛnā}, \textquotedblleft{}to drive away\textquotedblright{}.
\end{cousinweb}

\begin{grammarlens}[title={What you've built}]
One more word for getting around.
\end{grammarlens}

\subsection*{Your turn}
\begin{itemize}
\raggedright
  \item \emph{Say it:} \emph{khadeṛnā}
  \item \emph{Say it:} \emph{khadeṛnā}, once more
  \item \emph{Say it:} \emph{logoṁ ne bandroṁ ko galī se khadeṛā.} — People chased the monkeys out of the lane.
  \item \emph{Recall:} say the Hindi for a betel leaf, then the Hindi for an uproar, then say \emph{khadeṛnā} again
\end{itemize}

\begin{tcolorbox}[breakable,colback=teal!4,colframe=teal!35!black,title={Before you move on}]
What does \hi{खदेड़ना} mean? (\textquotedblleft{}To chase away\textquotedblright{}.) Say it once more.
\end{tcolorbox}
\section[\texorpdfstring{\hi{गटर} (gaṭar)}{गटर (gaṭar)}]{\hi{गटर} (gaṭar) — a sewer; a covered drain}

\label{lesson:HI-C425-gatar}



\begin{quote}
Before the new one: say the Hindi for an uproar, then the Hindi for to chase away.
\end{quote}

\subsection*{You'll want to know: \hi{गटर}}
\textbf{\hi{गटर}} — \emph{gaṭar} — \textquotedblleft{}a sewer; a covered drain\textquotedblright{}.

\begin{quote}
\hi{बारिश} \hi{में} \hi{गटर} \hi{भर} \hi{गया।}
\end{quote}

\begin{quote}
\emph{bāriś meṁ gaṭar bhar gayā.}
\end{quote}

\begin{quote}
The sewer filled up in the rain.
\end{quote}

Masculine.

\begin{cousinweb}
\textbf{\hi{गटर}} (\emph{gaṭar}) is borrowed from English \emph{gutter}; in Indian cities it names the covered sewer under the street.
\end{cousinweb}

\begin{grammarlens}[title={What you've built}]
One more word for getting around.
\end{grammarlens}

\subsection*{Your turn}
\begin{itemize}
\raggedright
  \item \emph{Say it:} \emph{gaṭar}
  \item \emph{Say it:} \emph{gaṭar}, once more
  \item \emph{Say it:} \emph{bāriś meṁ gaṭar bhar gayā.} — The sewer filled up in the rain.
  \item \emph{Recall:} say the Hindi for an uproar, then the Hindi for to chase away, then say \emph{gaṭar} again
\end{itemize}

\begin{tcolorbox}[breakable,colback=teal!4,colframe=teal!35!black,title={Before you move on}]
What does \hi{गटर} mean? (\textquotedblleft{}A sewer; a covered drain\textquotedblright{}.) Say it once more.
\end{tcolorbox}
\section[\texorpdfstring{\hi{बदबू} (badbū)}{बदबू (badbū)}]{\hi{बदबू} (badbū) — a bad smell, a stink}

\label{lesson:HI-C425-badbu}



\begin{quote}
Before the new one: say the Hindi for to chase away, then the Hindi for a sewer.
\end{quote}

\subsection*{You'll want to know: \hi{बदबू}}
\textbf{\hi{बदबू}} — \emph{badbū} — \textquotedblleft{}a bad smell, a stink\textquotedblright{}.

\begin{quote}
\hi{गटर} \hi{से} \hi{बहुत} \hi{बदबू} \hi{आ} \hi{रही} \hi{थी।}
\end{quote}

\begin{quote}
\emph{gaṭar se bahut badbū ā rahī thī.}
\end{quote}

\begin{quote}
There was a terrible smell coming from the sewer.
\end{quote}

Feminine.

\begin{cousinweb}
\textbf{\hi{बदबू}} (\emph{badbū}) is usually traced to Persian \emph{badbū}, from \emph{bad}, \textquotedblleft{}bad\textquotedblright{}, and \emph{bū}, \textquotedblleft{}smell\textquotedblright{}: the opposite of \hi{ख़ुशबू}.
\end{cousinweb}

\begin{grammarlens}[title={What you've built}]
One more word for getting around.
\end{grammarlens}

\subsection*{Your turn}
\begin{itemize}
\raggedright
  \item \emph{Say it:} \emph{badbū}
  \item \emph{Say it:} \emph{badbū}, once more
  \item \emph{Say it:} \emph{gaṭar se bahut badbū ā rahī thī.} — There was a terrible smell coming from the sewer.
  \item \emph{Recall:} say the Hindi for to chase away, then the Hindi for a sewer, then say \emph{badbū} again
\end{itemize}

\begin{tcolorbox}[breakable,colback=teal!4,colframe=teal!35!black,title={Before you move on}]
What does \hi{बदबू} mean? (\textquotedblleft{}A bad smell, a stink\textquotedblright{}.) Say it once more.
\end{tcolorbox}
\section[\texorpdfstring{\hi{रिसना} (risnā)}{रिसना (risnā)}]{\hi{रिसना} (risnā) — to seep; to leak}

\label{lesson:HI-C425-risna}



\begin{quote}
Before the new one: say the Hindi for a sewer, then the Hindi for a stink.
\end{quote}

\subsection*{You'll want to know: \hi{रिसना}}
\textbf{\hi{रिसना}} — \emph{risnā} — \textquotedblleft{}to seep; to leak\textquotedblright{}.

\begin{quote}
\hi{छत} \hi{से} \hi{पानी} \hi{रिस} \hi{रहा} \hi{था।}
\end{quote}

\begin{quote}
\emph{chat se pānī ris rahā thā.}
\end{quote}

\begin{quote}
Water was seeping through the roof.
\end{quote}

Past (it …): \hi{वह} \hi{तेल} \hi{डिब्बे} \hi{से} \hi{रिसा}, \emph{vah tel ḍibbe se risā}.

\begin{cousinweb}
\textbf{\hi{रिसना}} (\emph{risnā}) is filed by CDIAL under \emph{*riśyate}, \textquotedblleft{}is torn off\textquotedblright{}, a form it reconstructs (CDIAL 10745), with Hindi \emph{risnā}, \textquotedblleft{}to leak out, drip\textquotedblright{}.
\end{cousinweb}

\begin{grammarlens}[title={What you've built}]
One more word for getting around.
\end{grammarlens}

\subsection*{Your turn}
\begin{itemize}
\raggedright
  \item \emph{Say it:} \emph{risnā}
  \item \emph{Say it:} \emph{risnā}, once more
  \item \emph{Say it:} \emph{chat se pānī ris rahā thā.} — Water was seeping through the roof.
  \item \emph{Recall:} say the Hindi for a sewer, then the Hindi for a stink, then say \emph{risnā} again
\end{itemize}

\begin{tcolorbox}[breakable,colback=teal!4,colframe=teal!35!black,title={Before you move on}]
What does \hi{रिसना} mean? (\textquotedblleft{}To seep; to leak\textquotedblright{}.) Say it once more.
\end{tcolorbox}
